

En
composant (2.1.2) et (2.1.3), on obtient les homomorphismes (2.1.1).

La structure d'algèbre définie par les homomorphismes (2.1.1) est associative (cf. [1] XI 2). Par ailleurs, le \(A\)-module \(H_i(P_\bullet\otimes_A C)\) est muni de façon standard d'une structure de \(B\otimes_A C\)-module en faisant opérer \(C\) de façon évidente, et \(B\) par fonctorialité; de même, les \(H_j\bigl((P_\bullet\otimes_A C)\otimes_A(P_\bullet\otimes_A C)\bigr)\) sont munis de deux structures de \(B\otimes_A C\)-modules. Il est clair que les homomorphismes (2.1.2) et (2.1.3) sont compatibles avec ces deux structures: pour (2.1.2), par fonctorialité (en se plaçant dans la catégorie des \(A\)-modules), et pour (2.1.3), en utilisant le résultat classique sur les morphismes d'un complexe acyclique. Nous avons donc muni \(\bigoplus_{i\ge0}\Tor_i^A(B,C)\) d'une structure de \(B\otimes_A C\)-algèbre associative.

\subsection*{2.2.}
Supposons que \(B\) admette une résolution par des \(A\)-modules projectifs \(P_\bullet\), telle que \(P_\bullet=\bigoplus_{i\ge0}P_i\) soit une \(A\)-algèbre graduée, la différentielle satisfaisant:
\[
\begin{array}{ll}
\text{a)}& P_0\longrightarrow B\text{ est un homomorphisme d'anneaux~;}\SGAAuthorFootMark{1}\\[0.4ex]
\text{b)}& \forall x\in P_i,\ \forall y\in P_j,\quad d(xy)=d(x)y+(-1)^i x\,d(y).
\end{array}
\]
\SGAAuthorFootText{1}{On peut montrer qu'une telle résolution existe toujours.}
Alors les applications ``produit'' définissent des diagrammes commutatifs pour \(i,j\ge0\):
\begin{equation}\notag
\begin{tikzcd}[column sep=large,row sep=large]
P_i\otimes_A P_j \arrow[r] \arrow[d] & (P_{i-1}\otimes_A P_j)\oplus(P_i\otimes_A P_{j-1}) \arrow[d]\\
P_{i+j} \arrow[r] & P_{i+j-1},
\end{tikzcd}
\end{equation}
et par suite un homomorphisme de complexes \(P_\bullet\otimes_A P_\bullet\to P_\bullet\) qui prolonge l'homomorphisme d'augmentation : \(B\otimes_A B\to B\). Ceci montre que la structure d'algèbre de \(\bigoplus_{i\ge0}\Tor_i^A(B,C)\) peut se calculer à partir de celle de \(P_\bullet\). En particulier, si \(P_\bullet\) est une \(A\)-algèbre alternée (i.e. telle que pour tout
